% Guard: macros.tex is \input both by main.tex (after the packages it needs,
% e.g. amsmath for \DeclareMathOperator) and by individual section files so
% they remain compilable on their own. Loading it twice would redefine every
% \newcommand and abort, so the second and later loads are no-ops.
\ifdefined\PQRmacrosloaded\endinput\fi
\let\PQRmacrosloaded\relax

% --- 0. revision marking ---
% Text added to Section 2 during the current revision is wrapped in \added{...}
% so the advisors can see at a glance what is new. Redefine to \relax to print
% the paper without the marking.
\newcommand{\added}[1]{\textcolor{blue}{#1}}
\long\def\added#1{\begingroup\color{blue}#1\endgroup}  % \long: may span paragraphs


% --- A. operator-style score and function names ---
\newcommand{\HD}{\mathit{HD}}                        % histogram distance; use \HD_\gamma
\newcommand{\Attr}{\mathit{Attr}}                    % attribute set of a schema
\newcommand{\Attrs}{\mathit{Attrs}}                  % attribute set argument
% Abbreviated in MATH only; the prose still spells the names out in full.
% Spelling them out inside formulas makes displayed equations far too wide
% for a two-column layout (the product-rule step below overflowed by 200pt).
\newcommand{\Use}{\mathit{Use}}
\newcommand{\Cov}{\mathit{Cov}}
\newcommand{\Rep}{\mathit{Rep}}
\newcommand{\DivS}{\mathit{Div}}
\newcommand{\TVD}{\mathit{TVD}}
\newcommand{\GS}{\mathit{GS}}                        % global sensitivity
\DeclareMathOperator{\Dom}{Dom}
\DeclareMathOperator{\Lap}{Laplace}

% --- B. core objects ---
\newcommand{\Qone}{Q_1}
\newcommand{\Qtwo}{Q_2}
\newcommand{\pD}{p(D)}
\newcommand{\views}{\mathcal{P}}
\newcommand{\hist}[2]{H_{#1,#2}}                     % \hist{\Qone(\pD)}{A}
\newcommand{\dhist}{\Delta H}
\newcommand{\nhist}{\widetilde{\Delta H}}            % the released noisy object

% --- C. privacy budget and parameters ---
\newcommand{\eps}{\epsilon}
\newcommand{\epsP}{\epsilon_P}
\newcommand{\epsS}{\epsilon_S}
\newcommand{\epsH}{\epsilon_H}
\newcommand{\gammaref}{\gamma_{\mathrm{ref}}}
\newcommand{\gammaeval}{\gamma_{\mathrm{eval}}}
\newcommand{\taumax}{\tau_{\max}}

% --- D. mechanisms ---
\newcommand{\ExpMech}{\textsc{ExpMech}}
\newcommand{\OneShotTop}[1]{\textsc{OneShotTop-}#1}
\newcommand{\AdaptiveTau}{\textsc{AdaptiveTau}}
% --- E. experiment baselines (used in sections/experiments.tex) ---
\newcommand{\Bone}{\ensuremath{\mathsf{B1}}}
\newcommand{\Btwo}{\ensuremath{\mathsf{B2}}}
\newcommand{\Bthree}{\ensuremath{\mathsf{B3}}}
